\chapter{Conclusion}

\einlang{} is a tensor language built around three first-class constructs:
source-level multidimensional indexed tensor definitions, declarative
recurrences, and local derivative requests over named bindings. These constructs
form an executable abstraction boundary for tensor-heavy programs. The compiler
can use source-level axes, domains, dependency offsets, and derivative targets
to perform checks and analyses that are harder to recover after the program has
been split across host APIs.

The implementation supports this design with a full pipeline: Lark parsing,
source-located ASTs, DefId-based name resolution, module loading, IR lowering,
range and shape analysis, type inference, monomorphization, Einstein lowering,
recurrence order and storage metadata, compile-time and runtime autodiff,
validation, a DefId-keyed execution environment, a NumPy backend, an
experimental IREE backend, a standard library, examples, diagnostics, profiling,
and tests.

The next milestone is to deepen the implementation where the current thesis
draws conservative boundaries: mixed finite/full multidimensional storage
windows, richer dynamic-offset reasoning, wider IREE lowering, GPU execution,
and a stronger formal account of type and shape soundness. The architectural
through-line should remain the same: keep tensor, recurrence, and differential
structure visible in the source and preserve it until the compiler pass that can
use it.
